\section{GPU Modules: CUDA WPIC, Collisions, Restarts, and MPI}
\label{sec:gpu-modules}

\subsection{GPU-computing foundations}

A \emph{central processing unit} (CPU) is the general-purpose processor that
runs the operating system, reads files, chooses timesteps, and launches the
large arithmetic jobs in this chapter.  A \emph{graphics processing unit}
(GPU) is a second processor containing many simpler arithmetic workers.  It is
effective when the same operation must be applied independently to many
markers or mesh cells.  VENUS programs the GPU with NVIDIA's CUDA platform.

CUDA calls the CPU side the \emph{host} and the GPU side the \emph{device}.
Host and device normally have separate memory, so a \emph{memory transfer} is
an explicit copy between them.  A \emph{kernel} is a function launched on the
device.  One launch creates a \emph{grid} of \emph{blocks}; each block is a
group of \emph{threads}, and each thread is one logical worker executing the
kernel on its assigned marker, cell, or pair.  The CUDA launch grid organizes
workers; the radial mesh organizes the physical field cells. Threads in one
block can share a small fast memory area and coordinate with one another.

\emph{Synchronization} is a point at which selected workers must finish
before dependent work may continue.  A \emph{reduction} combines many values,
such as marker masses, into one sum or maximum.  An \emph{atomic update} is an
indivisible read--modify--write operation that prevents two threads from
losing each other's contributions when they update the same mesh cell.  These
mechanisms preserve contributions and the required operation dependencies.
The order of floating-point additions can still vary, giving CPU and GPU sums
different final rounding bits. \emph{Numerical parity} therefore requires
agreement within stated tolerances supported by the reduction accuracy,
together with conservation and convergence of the physical quantities.

The distinction between physics and acceleration is essential.  The Vlasov
equation, particle cloud, force solve, time integrator, and collision law
define the physical numerical algorithm.  Assigning their independent work
to CUDA threads, arranging device memory, and reducing partial results are
acceleration choices. Their implementation must retain the modeled equations,
boundary conditions, conservation laws and accuracy criteria while changing
how independent calculations are scheduled.

\noindent\textbf{Scope.}
The GPU code implements spherical weighted particle-in-cell
Vlasov--Poisson evolution, material SIDM collisions, calibrated subhalo
scattering and capture, versioned weighted-PIC restarts and an optional
one-MPI-rank-per-GPU layer. Einstein--Vlasov, the grid finite-volume and
semi-Lagrangian solvers, canonical-orbit Boltzmann--Poisson and SphDSMC remain
CPU modules. The following table locates the GPU calculations in the source.

\begin{center}
\small
\begin{tabular}{@{}p{0.38\linewidth}p{0.18\linewidth}p{0.34\linewidth}@{}}
\toprule
GPU component & Status & Primary source \\
\midrule
Persistent marker state and compaction & implemented &
\path{gpu/src/wpic_state.cu} \\
Spherical CIC/TSC Vlasov--Poisson KDK & implemented &
\path{gpu/src/wpic_vlasov_poisson.cu} \\
Material constant/power-law/Yukawa SIDM & implemented &
\path{gpu/src/wpic_collisions.cu} \\
Subhalo scatter/capture & implemented &
\path{gpu/src/wpic_subhalo_collisions.cu} \\
Restart versions 1--6 & implemented &
\path{gpu/src/wpic_restart.cpp} \\
One rank per GPU & implemented for VP and subhalo path &
\path{gpu/src/wpic_mpi.cpp} \\
Material-collision MPI invariance & not provided & rank partition changes the
fixed-occupancy shell schedule \\
GPU IC generation & not implemented & use a CPU weighted-PIC restart \\
\bottomrule
\end{tabular}
\end{center}

\subsection{State ownership and structure-of-arrays layout}

The public \code{venus_gpu::WpicState} owns persistent device arrays.  The
marker record is split into arrays so adjacent CUDA threads read adjacent
values:
\begin{itemize}
\item eight \code{double} arrays for represented mass, radius, radial
velocity, \(\eta=L^2/2\), inner-boundary speed and time, internal energy, and
internal velocity dispersion;
\item two \code{uint64_t} arrays for unique marker ID and source-cell ID;
\item one byte array for active/delayed/removed marker state.
\end{itemize}
Here \code{uint64_t} is an unsigned integer type with exactly 64 bits, used so
identifiers have a fixed nonnegative range across supported host and device
code. An ID labels a persistent marker and its random-number stream.
The arrays require 81 raw bytes per marker, supplemented by mesh workspaces,
sorting keys, scan storage, collision proposals and rollback buffers.
Initialization, requested output and checkpoints transfer the population.
Successive timesteps reuse its device arrays, reducing transfer cost and
preserving the ordered state used by the collision sampler.

CUDA 11.4 Thrust supports only ten fields in a zipped tuple.  Capture
compaction therefore uses two stable \code{copy_if} passes.  Stability is a
physics requirement here: retained marker IDs must remain in deterministic
order so a restart and the CPU pair scheduler see the same population.

\subsection{GPU Vlasov--Poisson equations}

For a spherical collisionless distribution in Newtonian gravity,
\begin{equation}
 \frac{\partial f}{\partial t}
 +w\frac{\partial f}{\partial r}
 +\left(\frac{L^2}{r^3}-\frac{\partial\Phi}{\partial r}\right)
  \frac{\partial f}{\partial w}=0,
 \qquad \eta=\frac{L^2}{2}.
\end{equation}
The field is
\begin{equation}
 \frac{1}{r^2}\frac{\partial}{\partial r}
 \left(r^2\frac{\partial\Phi}{\partial r}\right)=4\pi G\rho,
 \qquad
 \Phi_r(r)=\frac{GM(<r)}{r^2}.
\end{equation}
Here the distribution function determines the mass density, and the Poisson
potential determines radial acceleration. A collisionless marker keeps its
represented mass and \(\eta\) fixed while advancing \((r,w)\). Boundary
crossings update its boundary state; the collision operator below can also
remove a consumed capture partner.

\subsubsection{Spherical deposition}

The radial cloud calculation uses volume coordinate
\begin{equation}
 q=\frac{r^3}{3},\qquad
 \Delta V_i=4\pi(q_{i+1/2}-q_{i-1/2}).
\end{equation}
For q-CIC between centers \(i\) and \(i+1\),
\begin{equation}
 a_{pi}=\frac{q_{i+1}-q_p}{q_{i+1}-q_i},\qquad
 a_{p,i+1}=1-a_{pi}.
\end{equation}
Quantitative runs normally use the same volume-corrected q-TSC cloud as the
CPU implementation. The cloud spreads a marker's mass over nearby field
cells. The same weights deposit every extensive moment,
\begin{align}
 M_i &= \sum_p a_{pi}m_p,&
 P_i &= \sum_p a_{pi}m_pw_p,\\
 R_{2,i} &= \sum_p a_{pi}m_pw_p^2,&
 T_{2,i} &= \sum_p a_{pi}m_p\frac{2\eta_p}{r_p^2}.
\end{align}
The CUDA pattern is one thread per marker with atomic accumulation:
\par\noindent\begin{minipage}{\linewidth}
\codesource{GPU WPIC Vlasov--Poisson deposition}{\path{gpu/src/wpic_vlasov_poisson.cu}}{Assigns one marker to each CUDA thread and atomically deposits its spherical cloud moments into the radial mesh}
\begin{lstlisting}[language=C++]
const size_t p = blockIdx.x * blockDim.x + threadIdx.x;
if (p < marker_count && state[p] == active) {
  const double q = r[p] * r[p] * r[p] / 3.0;
  // Compute CPU-matched q-CIC or volume-corrected q-TSC weights.
  atomicAdd(&shell_mass[cell], cloud * mass[p]);
  atomicAdd(&radial_momentum[cell], cloud * mass[p] * w[p]);
}
\end{lstlisting}
\end{minipage}
Atomic arrival order can vary. Comparison with the CPU long-double sum
therefore checks discrete state, conservation residuals and stated rounding
tolerances for continuous moments.

An origin-parity carrier partitions a cloud that overlaps the first face.
For clipped smoothstep coordinate \(s\in[0,1]\),
\begin{equation}
 f_{\rm resolved}=s^2(3-2s),\qquad
 M_{\rm mesh}+M_{\rm hidden}=m_p.
\end{equation}
Each marker contributes its resolved fraction and complementary hidden
fraction. Their sum closes the marker's mass before the mesh totals are formed.

\subsubsection{Poisson scan and analytic potential}

For constant line mass in radial cell \(i\),
\begin{equation}
 \lambda_i=\frac{M_i}{\Delta r_i},\qquad
 M(r)=M_{i-1/2}+\lambda_i(r-r_{i-1/2}).
\end{equation}
Thrust is NVIDIA's C++ library of parallel algorithms.  Its
\emph{inclusive scan} is a prefix sum: output entry \(i\) contains the sum of
all input entries from zero through \(i\), including entry \(i\) itself.
Applied to shell masses, it constructs the enclosed mass at successive radial
edges.  Independent cell threads then evaluate the center force.  With
\(A_i=M_{i-1/2}-\lambda_i r_{i-1/2}\), the exact cell contribution to the
potential is
\begin{equation}
 \int_{r_a}^{r_b}\frac{GM(r)}{r^2}\,\dd r
 =G\left[A_i\left(\frac1{r_a}-\frac1{r_b}\right)
 +\lambda_i\ln\!\left(\frac{r_b}{r_a}\right)\right].
\end{equation}
A reverse scan begins from
\(\Phi(r_{\max})=-GM(r_{\max})/r_{\max}\). For total-density input in an
expanding background, the source subtracts the homogeneous density
\(\bar\rho\) in both the resolved shells and the unresolved central sphere.
For density-contrast input, the marker masses already represent the source.

The scan begins at the inner mesh face, \(r_{\min}\), with the sum of
unresolved marker mass and any prescribed central point mass. The
total-density background correction there is
\(-4\pi\bar\rho r_{\min}^3/3\). Below that face the unresolved distributed
source is uniform, as in (PIC11a)--(PIC11b). The face mass fixes its force;
matter in the first resolved shell contributes through the exterior field.
\codesource{GPU force inside the first radial face}
{\path{gpu/src/wpic_vlasov_poisson.cu}}
{Use the enclosed face mass to continue the distributed central source}
\begin{lstlisting}[language=C++]
const double inner = edge[0];
const double core_coefficient = options.gravity
    * (source_mass_edge[0] - options.central_mass)
    / (inner * inner * inner);
\end{lstlisting}
The potential uses the same face anchor, so both potential and force are
continuous across the face. Refining the inner mesh tests the uniform-source
approximation where central structure is scientifically relevant.

\subsubsection{KDK and exact central drift}

The first half-kick in a static background is
\begin{equation}
 w^{n+1/2}=w^n-\frac{\Delta t}{2}\Phi_r^n(r^n).
\end{equation}
The drift integrates free central motion analytically, including the changing
radial component of the tangential velocity.  With \(L=\sqrt{2\eta}\),
\begin{align}
 r_*^2&=(r^n+\Delta t\,w^{n+1/2})^2
       +\left(\Delta t\frac{L}{r^n}\right)^2,\\
 w_*&=\frac{w^{n+1/2}(r^n+\Delta t\,w^{n+1/2})
 +(L/r^n)(\Delta t L/r^n)}{r_*}.
\end{align}
After endpoint deposition and a second Poisson solve,
\begin{equation}
 w^{n+1}=w_*-\frac{\Delta t}{2}\Phi_r^{n+1}(r_*).
\end{equation}
For \(L>0\), the implementation keeps canonical momentum in an expanding background:
\par\noindent\begin{minipage}{\linewidth}
\codesource{GPU WPIC orbit pusher}{\path{gpu/src/wpic_vlasov_poisson.cu}}{Applies the expanding-background half kick and exact central drift to each active marker}
\begin{lstlisting}[language=C++]
double p = a0 * w - 0.5 * dt * a0 * phi_r;
const double radial_axis = x + drift * p;
const double tangential_axis = drift * L / x;
const double new_x = hypot(radial_axis, tangential_axis);
p = (p * radial_axis + (L / x) * tangential_axis) / new_x;
r = a1 * new_x;
w = p / a1;
\end{lstlisting}
\end{minipage}
For \(r=a(t)x\) and \(p=aw\), the drift factor is
\begin{equation}
 S(t_n,t_{n+1})=\int_{t_n}^{t_{n+1}}\frac{\dd t}{a(t)^2}.
\end{equation}
For \(a=[(t+t_0)/t_0]^\alpha\), the code evaluates this power-law integral
analytically and uses its logarithmic limit near \(1-2\alpha=0\).  The GPU
option preflight enforces the same expansion domain as the CPU solver:
\(\alpha\geq0\), \(t_0>0\), and \(t+t_0>0\), with every participating scalar
finite.  Homogeneous-background subtraction additionally requires \(G>0\);
ballistic \(G=0\) evolution remains valid only when the input marker masses
already represent the desired contrast source and subtraction is disabled.

A reflecting wall is applied at the time when this free trajectory reaches
its fixed physical radius \(R\). The normal physical velocity changes sign,
so \(w_{\rm out}=-w_{\rm in}-2H(t_{\rm hit})R\), while angular momentum
remains constant. The shared CPU/CUDA routine in
\path{cpu/include/ev/central_drift.hpp} locates the impact by dividing the
trajectory at its turning points and bisecting each monotonic interval.
After reflection it advances the unused part of the drift. This resolves
inner-wall impacts even when the unconstrained orbit returns outside the
wall before the full step ends.

For exactly radial motion, \(L=0\), the signed comoving coordinate advances
as \(x+p\int a^{-2}\dd t\). Its absolute value gives the new comoving
radius. Crossing zero reverses the radial momentum. A forward step ending
exactly at the origin stores the outgoing velocity, so the next step
continues the orbit. The implementation treats this case directly:
\codesource{Radial passage through the origin}
{\path{gpu/src/wpic_vlasov_poisson.cu}}
{Advance the signed radial coordinate and continue the orbit through zero}
\begin{lstlisting}[language=C++]
const double signed_x = x + drift * p;
if (signed_x < 0.0) p = -p;
else if (signed_x == 0.0 && dt > 0.0) p = fabs(p);
r = a1 * fabs(signed_x);
w = p / a1;
\end{lstlisting}
An origin start requires \(L=0\) and zero central point mass. The smooth
central force then vanishes at the origin. With a nonzero central point mass,
the force is singular there, and the solver rejects a radial step that
reaches or crosses the origin. CPU and GPU use the same convention. The
shared wall solver also accepts an origin start followed by a finite
outer-wall impact within the same drift.

The collisionless task uses the fixed KDK orbit step
\code{--maximum-timestep}; only an output or final-time landing may shorten
it.  This preserves the time symmetry and bounded energy error of KDK in a
constant potential.  The cell dynamical estimate is used only by collisional
tasks, whose event splitting already breaks the collisionless fixed-step
symmetry.

\subsection{GPU material SIDM collision module}

The material and subhalo operators solve
\begin{equation}
 \frac{\partial f}{\partial t}+\{f,H\}=C[f]
\end{equation}
by sampling local particle pairs. Each family specifies its physical collision
law and the probability with which a pair is proposed.  For material scattering let
\(u=|\bm v_i-\bm v_j|\).  The configured law is area per represented mass,
\(s(u)=\sigma(u)/m\), and the unequal-marker closure is
\begin{equation}
 \Sigma_{ij}(u)=\frac{m_i+m_j}{2}s(u).
\end{equation}
The GPU supports
\begin{align}
 s_{\rm const}(u)&=S\sigma_0,\\
 s_{\rm pow}(u)&=S\sigma_0(u_*^2+u_0^2)^{-n/2},
 \qquad u_*=\max(|u|,10^{-14}),\\
 s_{\rm Y}(u)&=
 S\frac{4\pi\alpha_\chi^2m_\chi^2}
 {m_\phi^4[1+m_\chi^2u_*^2/m_\phi^2]}.
\end{align}
The Yukawa--Born differential law per represented mass is
\begin{equation}
 \frac{\dd s}{\dd\Omega}=
 S\frac{\alpha_\chi^2m_\chi^2}
 {[m_\phi^2+\tfrac12m_\chi^2u^2(1-\mu)]^2},
 \qquad \mu=\cos\theta.
\end{equation}
The numerical parameters set the amplitude and speed scale after the physical
conversion in Section~\ref{sec:yukawa-material-units}.  The GPU uses that same
material area-per-mass convention.
The angular inverse uses \(y=(m_\chi u/m_\phi)^2\) and a uniform
upper-tail probability \(z\in(0,1)\):
\begin{equation}
 \mu=1-\frac{2z}{1+y(1-z)}.
\end{equation}
This approaches the uniform cosine \(1-2z\) as \(y\to0\).  For
\(y>1\), the implementation uses the equivalent expression with numerator
and denominator divided by \(y\).

Active markers are sorted by \((r,{\rm id})\) into target-occupancy shells.
A group with equal radii stays in one shell: the boundary advances through
the whole group, and a final underfilled shell is merged. This gives the
collision sampler positive shell volumes even for cell-centred particles.
One block owns a shell; shared reductions construct a conservative majorant.
For a shell with population \(N_s\ge2\) and volume \(V_s\), the pair
scheduler uses
\begin{equation}
 m_s=\begin{cases}N_s-1,&N_s\ \hbox{even},\\
                   N_s,&N_s\ \hbox{odd}.
       \end{cases}
\end{equation}
These \(m_s\) partner-assignment rounds contain every unordered pair once.
Each round assigns every marker at most one partner; for odd \(N_s\), one
marker sits out.  A uniformly drawn round therefore includes each pair with
probability \(1/m_s\).  A rejection draw from a 64-bit random word makes the
round selection uniform.  The disjoint assignments let separate CUDA
threads update their two markers safely:
\par\noindent\begin{minipage}{\linewidth}
\codesource{GPU material-SIDM pair scheduler}{\path{gpu/src/wpic_collisions.cu}; \path{gpu/include/venus_gpu/collision_matching.hpp}}{Maps the shared random round and thread-owned pair slot to two distinct markers}
\begin{lstlisting}[language=C++]
const collision_matching::Pair proposed =
    collision_matching::pair_at(count, matching_round, pair);
const std::uint64_t first_local = proposed.first;
const std::uint64_t second_local = proposed.second;
const std::uint64_t first_index = order[shell.begin + first_local];
const std::uint64_t second_index = order[shell.begin + second_local];
\end{lstlisting}
\end{minipage}
For an assigned pair, the acceptance probability over local step
\(\delta t_s\) is
\begin{equation}
 P_{ij}=\frac{m_s}{V_s}u\Sigma_{ij}(u)\,\delta t_s.
\end{equation}
Multiplying by the inclusion probability \(1/m_s\) gives the physical
pair probability \(u\Sigma_{ij}\delta t_s/V_s\).  The physical partner
number density \((N_s-1)/V_s\) continues to set the mean-free-path diagnostic;
the proposal exposure \(m_s/V_s\) sets acceptance and its timestep bound.
Let \(B_s\) be the safe shell bound on \(u\Sigma_{ij}\), constructed from
the maximum marker mass and speed.  The collision rate bound is
\begin{equation}
 \nu_s=\frac{m_s}{V_s}B_s,\qquad
 B_s\ge\max_{ij\in s}(u\Sigma_{ij}).
\end{equation}
Reference mode advances every occupied shell with the common step
\begin{equation}
 \delta t_{\rm fixed}=\min\left(
 \Delta t_{\rm batch}-t,
 \frac{C_{\rm coll}}{\max_s\nu_s}\right).
\end{equation}
With \code{--sidm-adaptive-shell-timestep}, each shell instead stores \(t_s\)
and takes
\begin{equation}
 \delta t_s=\min\left(
 \Delta t_{\rm batch}-t_s,
 \frac{C_{\rm coll}}{\nu_s}\right).
\end{equation}
All currently active shells share a CUDA launch.  The shell-description kernel
then refreshes maximum speed and the safe \(u\Sigma\) majorant before the next
local step.  Elastic scattering freezes radii and masses, hence shell
membership, faces, density, and volume remain valid.  This device path supports
constant, power-law, and Yukawa--Born laws and unequal marker masses.  In both
modes \(\nu_s\delta t_s\le C_{\rm coll}\le1\); local clocks only avoid
revisiting slow shells at every dense-core step.
The host/device schedule is structurally
\par\noindent\begin{minipage}{\linewidth}
\codesource{GPU material-SIDM adaptive shell clock}{\path{gpu/src/wpic_collisions.cu}}{Advances independent host-owned shell clocks and launches the device collision and shell-refresh kernels}
\begin{lstlisting}[language=C++]
while (some_shell_is_unfinished) {
  for (shell : shells)
    shell_dt[shell] = min(batch_dt - elapsed[shell],
                          collision_cfl / rate[shell]);
  upload(shell_dt);
  collide_shells_kernel<<<number_of_shells, threads>>>(...);
  describe_shells_kernel<<<number_of_shells, threads>>>(...);
}
\end{lstlisting}
\end{minipage}
The pair marginals reproduce the instantaneous Boltzmann collision rate.
Each marker experiences at most one collision per local step, so quantitative
evolution requires convergence as \(C_{\rm coll}\) is reduced.  Material
CPU/GPU comparison checks formulas, conservation, timestep convergence, and
statistical ensembles.  MPI rank ownership changes the fixed-occupancy
shells, so multi-GPU calculations also require rank-count convergence.

\subsection{Unequal-mass elastic collision map}

Both GPU collision families lift \((w,\eta)\) to a three-vector using a
counter-RNG azimuth.  Define
\begin{equation}
 \bm V=\frac{m_i\bm v_i+m_j\bm v_j}{m_i+m_j},\qquad
 \bm u=\bm v_i-\bm v_j.
\end{equation}
After drawing \(\widehat{\bm n}\) and setting
\(\bm u'=u\widehat{\bm n}\),
\begin{align}
 \bm v_i'&=\bm V+\frac{m_j}{m_i+m_j}\bm u',\\
 \bm v_j'&=\bm V-\frac{m_i}{m_i+m_j}\bm u'.
\end{align}
Projection back uses \(w'=v_x'\) and
\(\eta'=r^2(v_y'^2+v_z'^2)/2\).  Pair momentum and kinetic energy are
conserved at roundoff even for unequal masses.

\subsection{GPU subhalo scatter and capture}

The build converts the maintained CPU Plummer table into a CUDA header.
Tracking that input as a build dependency gives both evaluators the same
calibration whenever the table changes.  For masses \(M_i,M_j\),
\begin{equation}
 A_{ij}=\frac{A}{4}(\sqrt{M_i}+\sqrt{M_j})^2,
 \qquad q=\frac{\min(M_i,M_j)}{\max(M_i,M_j)},
\end{equation}
and
\begin{equation}
 \sigma_{ij}^2=
 \frac{\sqrt{M_i}\sigma_i^2+\sqrt{M_j}\sigma_j^2}
 {\sqrt{M_i}+\sqrt{M_j}},
 \qquad x=\frac{u}{\sigma_{ij}}.
\end{equation}
The mutually exclusive cross sections are
\begin{equation}
 \Sigma_{\rm scatt}=SA_{ij}F_{\rm scatt}(x,q),\qquad
 \Sigma_{\rm cap}=SA_{ij}F_{\rm cap}(x,q).
\end{equation}
The table evaluator uses log-speed/log-mass-ratio interpolation, zero-aware
closure, low-speed regularization, vanishing high-speed capture, and an
\(x^{-4}\) scattering continuation. The pair area sets the dimensional
normalization; the functions of speed and mass ratio supply the calibrated
scattering and capture probabilities.

This module uses the CPU
\code{maximum_weight_importance} pair-majorant schedule. A pair majorant is
an upper bound on the physical pair rate, used to propose candidate events.  Within each
physical radial cell it sorts by internal dispersion and ID.  For proposal
positions \(i<j\), a safe weight has form
\begin{equation}
 B_{ij}=C\,\sigma_j(\sqrt{M_i}+\sqrt{M_j})^2.
\end{equation}
Prefix sums of \(\sqrt M\) and \(M\) sample the complete pair distribution in
\(O(N\log N)\) work and \(O(N)\) storage.  The stochastic trial count and
acceptance probability are
\begin{equation}
 N_{\rm trial}=\left\lfloor
 \frac{\Delta t_c}{V}\sum_{i<j}B_{ij}+z_{\rm round}
 \right\rfloor,
 \qquad
 p_{ij}=\frac{u(\Sigma_{\rm scatt}+\Sigma_{\rm cap})}{B_{ij}}\le1.
\end{equation}
One accepted event per marker per substep is enforced.  Proposal arrays remain
cached after elastic scattering because their masses and internal dispersions
did not change.  A capture invalidates the cache, compacts the population, and
rebuilds proposals before the captured cell advances again.

Let \(\nu_{B,c}=\max_{k\in c}\nu_{B,k}\).  Reference mode uses
\begin{equation}
 \delta t_{\rm fixed}=\min\left(
 \Delta t_{\rm batch}-t,
 \frac{C_{\rm coll}}{\max_c\nu_{B,c}}\right)
\end{equation}
for every populated cell.  With \code{--subhalo-adaptive-cell-timestep}, cell
\(c\) owns \(t_c\) and takes
\begin{equation}
 \delta t_c=\min\left(
 \Delta t_{\rm batch}-t_c,
 \frac{C_{\rm coll}}{\nu_{B,c}}\right).
\end{equation}
One CUDA thread owns a cell and performs at most
\code{--subhalo-adaptive-batch-substeps} local steps per kernel launch.
Elastic encounters reuse the proposal.  A capture terminates that cell's
device loop immediately; host orchestration compacts the population, rebuilds
proposals, and resumes unfinished cells from their saved elapsed time and
local substep index.  \code{maximum_cell_substeps}, \code{cell_visits},
\code{kernel_batches}, and \code{proposal_rebuilds} expose this scheduling
work without changing the physical event counters.
The device-owned clock is structurally
\par\noindent\begin{minipage}{\linewidth}
\codesource{GPU subhalo adaptive cell clock}{\path{gpu/src/wpic_subhalo_collisions.cu}}{Lets one CUDA thread advance one collision cell through bounded local substeps and stop immediately after capture}
\begin{lstlisting}[language=C++]
while (cell_elapsed[cell] < batch_dt && local_steps < batch_limit) {
  dt_cell = min(batch_dt - cell_elapsed[cell],
                collision_cfl / cell.maximum_trial_frequency);
  captured = collide_one_cell_substep(cell, dt_cell, local_substep);
  cell_elapsed[cell] += dt_cell;
  if (captured) break;  // compact and rebuild before resuming
}
\end{lstlisting}
\end{minipage}

Capture merges complete objects.  For \(M=M_i+M_j\) and reduced mass
\(\mu_{\rm red}=M_iM_j/M\),
\begin{align}
 \bm v_{\rm rem}&=\bm V,
 &r_{\rm rem}&=\frac{M_ir_i+M_jr_j}{M},\\
 E_{{\rm int},{\rm rem}}&=E_{{\rm int},i}+E_{{\rm int},j}
 +\frac12\mu_{\rm red}u^2,
 &\sigma_{\rm rem}&=\sqrt{\frac{2E_{{\rm int},{\rm rem}}}{3M}}.
\end{align}
Thus relative orbital kinetic energy becomes internal heat while total mass,
vector momentum, and pair kinetic-plus-internal energy close.  Moving the
deposited mass to the remnant radius can change the reconstructed potential
energy; the complete gravitational budget and collision-cell refinement
measure that contribution.

The locality measure is
\begin{equation}
 \lambda_{\rm mfp}=
 \frac{\langle u\rangle}
 {n_{\rm partner}\langle u\Sigma_{ij}(u)\rangle}.
\end{equation}
The runner records \(\max(\Delta r/\lambda_{\rm mfp})\), warns when it
exceeds unity, and stops when it exceeds the configured hard limit.  Collision coverage should also be reported using
\begin{equation}
 \tau_{\rm coll}(t)=\int^t\frac{\dd t'}{t_{\rm coll}(t')},
\end{equation}
alongside the final physical time. A history column containing the minimum
cell collision time measures the fastest local rate. A halo-average collision
age requires the chosen population average integrated through its evolution.

\subsection{One MPI rank per GPU}

The device-binding helper forms a node-local communicator and maps each local
rank to one visible device.  For Vlasov--Poisson, rank \(a\) deposits local
moments \(M_i^{(a)}\); host \code{MPI_Allreduce} constructs
\begin{equation}
 M_i=\sum_a M_i^{(a)}
\end{equation}
and the other \(O(N_r)\) moment arrays.  Every GPU uploads the same global
moments, solves the same Poisson problem, and pushes only its local markers.
KDK therefore performs two radial-moment reductions per step while each
rank retains its own markers.

The local collision sampler requires complete physical radial cells. These
cells are assigned to contiguous rank slabs; markers migrate through host
records with \code{MPI_Alltoallv}.  Random draws retain the global tuple
\par\noindent\begin{minipage}{\linewidth}
\codesource{Multi-GPU subhalo collision sampler}{\path{gpu/src/wpic_mpi.cpp} and \path{gpu/src/wpic_subhalo_collisions.cu}}{Assigns complete radial cells to ranks and keys every counter-based random draw by a rank-independent global collision identity}
\begin{lstlisting}
(seed, epoch, global collision cell, substep, trial, dimension)
\end{lstlisting}
\end{minipage}
Fixed reference mode uses \code{MPI_Allreduce(MPI_MAX)} to obtain its common
mesh-wide collision step. In adaptive mode the cells already have unique
owners, so each advances independently inside the synchronized outer operator
batch.  Global cell numbers preserve the counter-RNG definition.
The migration supplies those complete populations and may dominate a short
run with frequent collision calls.

At checkpoint time, rank zero gathers complete marker records, sorts them by
unique ID, and writes the ordinary version-6 checkpoint.  Restart clocks,
RNG epoch, step, time, collision age, and potential-event age are preserved.

\subsection{Optional receiver field clocks for subhalo collisions}

The single-GPU subhalo task accepts \code{--coupled-receiver-clocks}
with adaptive collision cells and a positive collision-timescale fraction.
The orbit bound uses the existing positive-source cell floor and
\begin{equation}
 \Delta t_{\rm grav}=\min\left[\Delta t_{\max},
 f_{\rm dyn}\min_i\left(|\Phi_{r,i}|/r_i+H_*^2\right)^{-1/2}\right],
 \qquad H_*={2\over3(t+12)}.
\end{equation}
The parameter \(f_{\rm dyn}\) is the dynamical fraction. The minimum
includes source cells above \(10^{-12}\) of the deposited source mass.
The numerical frequency envelope retains \(H_*\) for a static background;
its physical expansion rate and cosmological acceleration remain zero.
All represented matter remains in the force source and characteristic update.

For positive gravity, an already background-subtracted EdS source, TSC
deposition, retained inner mass and zero central mass, the field
observer uses resolved deposited shell masses and the positive hidden
cloud masses inside the first face. For either potential or force
values \(A\), let \(\omega_k\) denote receiver mass and let \(A_{0,k}\)
and \(A_{1,k}\) be the fields before and after an observation interval
\(h\). Their relative change is
\begin{equation}
 \delta_A^2={\sum_k\omega_k(A_{1,k}-A_{0,k})^2\over
 \max(\sum_k\omega_k A_{0,k}^2,\sum_k\omega_k A_{1,k}^2)}.
\end{equation}
The resolved deposition supplies
\(N_{\rm eff}=(\sum_i M_i)^2/\sum_i M_i^2/N_{{\rm eff},i}\),
where \(M_i\) and \(N_{{\rm eff},i}\) are resolved shell masses and
effective marker counts. The sampling floor is
\(\epsilon_{\rm samp}=\max(4096\epsilon_{\rm fp},0.25/\sqrt{N_{\rm eff}})\), where
\(\epsilon_{\rm fp}\) is binary64 machine precision. Above this floor,
the selected channel time is \(h/\sqrt{\delta_A^2-\epsilon_{\rm samp}^2}\).
At or below the floor the channel time is infinite. The smaller potential/force time
controls the field contribution. Static and other unsupported observer
configurations use the fallback maximum-norm field clock. Interior
receivers use the uniform-density inner continuation at particle radii.
Their hidden cloud fractions follow deposition's separately rounded
operations. Hidden mass closes within \(2048\epsilon_{\rm fp}\) relative
error, and resolved plus hidden receiver mass closes the source ledger
within \(2\times10^{-11}\). A failure preserves diagnostics and stops.

The additional \code{--collision-prestep-coordinate-due} option applies
a positive carried collision age \(a\) at the current particle state when
its fresh deadline \(d\) satisfies
\begin{equation}
 a>0,\qquad d>q,\qquad d-a\le q,\qquad
 q=128\epsilon_{\rm fp}\max(1,|t|).
\end{equation}
The age must also satisfy the existing overdue screen
\(a\le d+64\epsilon_{\rm fp}\max(1,|a|,|d|)\). The unchanged sampler
consumes that age and advances its persistent random counter.
Trial endpoints still recompute the field and collision deadline before
commitment, and unsafe trials restore the incoming particle state before
retry. Locality, collision exposure and capacity checks retain their
existing thresholds. The policy defaults to off and requires a fresh output
segment with the complete carried restart clocks. The source equations,
supported backgrounds and diagnostic records are described in
\path{gpu/docs/COUPLED_RECEIVER_CLOCKS.md}.
The optional \code{--coupled-field-trace} writes receiver and cadence
observations. It requires \code{--coupled-receiver-clocks}; its observations
preserve physical state and random counters. The CPU expanding WPIC driver
exposes the same default-off pre-step due option through the shared cadence
helper.

\subsection{GPU build and run}

From the \path{VENUS} root, the complete single-GPU build is
\par\noindent\begin{minipage}{\linewidth}
\codesource{Single-GPU build}{\path{gpu/CMakeLists.txt}}{Configures the CUDA libraries, task executables, benchmarks, generated calibration table, and shared production drivers}
\begin{lstlisting}[language=bash]
cmake -S gpu -B build/gpu \
  -DCMAKE_BUILD_TYPE=Release \
  -DCMAKE_CUDA_ARCHITECTURES=80 \
  -DVENUS_CPU_REFERENCE_INCLUDE="$PWD/cpu/include"
cmake --build build/gpu -j 8
\end{lstlisting}
\end{minipage}
The bundled CPU include path is selected automatically when the trees remain
siblings.  A full build contains all six task-specific runners, the restart
comparator, focused VP/material benchmarks, and the generated subhalo table.
The host-side comparator can be used without a visible GPU after the build.

MPI adds
\par\noindent\begin{minipage}{\linewidth}
\codesource{One-rank-per-GPU build}{\path{gpu/CMakeLists.txt}}{Builds the MPI-enabled CUDA layer and exposes its production interface}
\begin{lstlisting}[language=bash]
cmake -S gpu -B build/gpu-mpi \
  -DVENUS_GPU_ENABLE_MPI=ON \
  -DVENUS_CPU_REFERENCE_INCLUDE="$PWD/cpu/include" \
  -DCMAKE_CUDA_ARCHITECTURES=80
cmake --build build/gpu-mpi -j 8
srun --ntasks=2 --gpus-per-task=1 build/gpu-mpi/venus_gpu_vlasov_poisson_mpi --help
\end{lstlisting}
\end{minipage}
That configuration also builds
\path{gpu/benchmarks/benchmark_wpic_mpi.cu}, the global one-rank-per-GPU
throughput harness.
Use one rank per GPU and no more ranks than radial cells.

Every GPU runner requires a weighted-PIC restart.  The subhalo task is:
\par\noindent\begin{minipage}{\linewidth}
\codesource{GPU subhalo production task}{\path{gpu/tools/venus_gpu_subhalo_boltzmann_poisson.cu} and \path{gpu/src/wpic_production_driver.cu}}{Selects the subhalo model and runs its shared restart, option, evolution, collision, and output driver}
\begin{lstlisting}[language=bash]
build/gpu/venus_gpu_subhalo_boltzmann_poisson \
  --restart restart_state.bin \
  --output-directory output/gpu/run \
  --end-time 60 \
  --subhalo-collision-strength 0.1
\end{lstlisting}
\end{minipage}
Use \code{venus_gpu_vlasov_poisson} for the collisionless task and
\code{venus_gpu_sidm_boltzmann_poisson}, \code{venus_gpu_sidm_constant},
\code{venus_gpu_sidm_power_law}, or \code{venus_gpu_sidm_yukawa_born} for
static material SIDM.  Appending \code{_mpi} selects one rank per GPU.
All tasks fix \(G=1\), TSC deposition, origin parity and a reflecting
outer wall. Single-GPU VP and subhalo tasks accept static total-density evolution or
Einstein--de Sitter density-contrast evolution, with \(\alpha=2/3\) and time
offset 12. EdS is the VP default. MPI builds restrict VP and subhalo tasks to EdS.
Material tasks use a static, total-density halo. The VP fixed KDK cap and
output/final-time clipping apply in both admitted single-GPU backgrounds.  Initial conditions,
including \(\gamma_i\), come from the restart; GPU runners do not generate
them.

Table~\ref{tab:gpu-task-cli} gives the common and subhalo
command-line contract.  Defaults in brackets apply when an optional switch is
omitted.  A task executable rejects options from another collision family.

\small
\begin{longtable}{@{}p{0.43\linewidth}p{0.17\linewidth}p{0.34\linewidth}@{}}
\caption{GPU common and subhalo command-line interface.}
\label{tab:gpu-task-cli}\\
\toprule
Option & Default or admissible range & Physical or numerical meaning \\
\midrule
\endfirsthead
\multicolumn{3}{l}{\small\itshape Table~\ref{tab:gpu-task-cli} continued}\\
\toprule
Option & Default or admissible range & Physical or numerical meaning \\
\midrule
\endhead
\midrule
\multicolumn{3}{r}{\small\itshape Continued on next page}\\
\endfoot
\bottomrule
\endlastfoot
\path{--restart FILE} & required & Versioned CPU/GPU WPIC checkpoint.  All
imported markers must currently be active, and the requested end time must
follow the checkpoint time.  A subhalo collisional run (\(S>0\))
also requires a positive internal velocity dispersion for every participating marker and
the stored closure
\(E_{{\rm int},p}=3M_p\sigma_{s,p}^2/2\) to roundoff. \\
\path{--output-directory DIR} & required & Run directory in which profiles,
histories, metadata, checkpoints, and terminal status are published. \\
\path{--end-time T} & required; \(T>0\) and \(T>t_{\rm restart}\) & Absolute
final simulation time, not a duration measured from the restart. \\
\path{--universe MODE} & VP default \code{eds} & Single-GPU VP accepts
\code{static} or \code{eds}. MPI VP and subhalo tasks use \code{eds};
material SIDM uses \code{static}. \\
\path{--maximum-timestep DT} & \(0.015;\ DT>0\) & Fixed collisionless KDK
step, apart from output/final landings; collisional tasks treat it as the
absolute orbit-step ceiling before event deadlines. \\
\path{--dynamical-fraction F} & \(0.04;\ F>0\) & Multiplier of the most
restrictive cell dynamical estimate used by the orbit-step controller. \\
\path{--output-interval DT} & \(10;\ DT>0\) & Simulation-time separation of
radial profile outputs on an absolute cadence anchored at \(t=0\). \\
\path{--history-every N} & \(5000;\ N\ge1\) integer & Append and flush the
scalar evolution history every \(N\) accepted KDK steps; an output or final
step also forces a row. \\
\path{--restart-every N} & \(25000;\ N\ge1\) integer & Write an atomic
versioned checkpoint every \(N\) accepted KDK steps; final and controlled-stop
paths also checkpoint. \\
\path{--radial-cells N} & \(256;\ N\ge2\) integer & Number of cells in the
fixed logarithmic field mesh. \\
\path{--radial-minimum R} & \(0.1;\ R>0\) & Inner face of the logarithmic field
mesh.  Active mass inside it is carried by the origin-parity treatment. \\
\path{--radial-maximum R} & \(1000;\ R>R_{\min}\) & Reflecting outer face of
the logarithmic field mesh. \\
\path{--subhalo-collision-strength S} & required; \(S>0\) & Dimensionless
multiplier \(S\) of both calibrated subhalo scatter and capture cross
sections.  This option is valid only for the subhalo task. \\
\path{--subhalo-collision-interval DT} & \(0.5;\ DT>0\) & Hard upper bound on
the physical time between subhalo collision calls; it is the fixed
interval when the adaptive fraction is zero. \\
\path{--subhalo-collision-timescale-fraction F} & \(0.05;\ 0\le F\le1\) &
Adaptive fraction of the controlling collision, dynamical, and potential
timescales.  \(F=0\) selects only the fixed maximum interval. \\
\path{--subhalo-collision-cfl C} & \(0.2;\ 0<C\le1\) & Bounds the maximum
trial probability per collision substep and therefore controls stochastic
time subdivision. \\
\path{--subhalo-collision-maximum-substeps N} & \(16384;\ N\ge1\) integer &
Administrative safety ceiling on substeps in one collision call; exceeding it
is an error, not permission to under-resolve the event. \\
\path{--subhalo-adaptive-cell-timestep} & off & Give every collision cell an
independent CFL clock inside one synchronized operator batch. \\
\path{--subhalo-adaptive-batch-substeps N} & \(256;\ N\ge1\) integer & Maximum
number of local cell steps performed by one CUDA thread in a kernel launch;
capture always stops the cell earlier for proposal reconstruction. \\
\path{--subhalo-inner-collision-cells N} & \(1;\ 1\le N\le256\) integer &
Number of equal-volume virtual collision cells inside the field mesh's inner
face.  These resolve carried inner mass without changing the field mesh. \\
\path{--subhalo-collision-lanes N} & \(8;\ 1\le N\le256\) integer & Number of
deterministic ID/epoch strata per physical collision cell.  Lanes partition
pairing work and cell volume; they are not extra radial field cells. \\
\path{--subhalo-locality-maximum-ratio X} & \(10;\ X\ge0\) & Controlled-stop
threshold for \(\max(\Delta r/\lambda_{\rm mfp})\).  \(X=0\) disables the hard
stop, but the diagnostic warning above unity remains active. \\
\path{--subhalo-velocity-scale V} & \(0.1;\ V>0\) & Compatibility control,
retained and range-checked by the interface.  This restart-only runner instead
derives every pair scale from the markers' stored internal velocity
dispersions, so changing \(V\) alone does not rescale the current collision
kernel. \\
\path{--subhalo-plummer-area-per-mass A} & \(1;\ A>0\) & Plummer geometric
area per mass in simulation units, entering \(A_{ij}\) in the cross-section
normalization. \\
\path{--subhalo-low-speed-regularization E} & \(0.05;\ 0.05\le E<1\) & Lower
dimensionless-speed regularization for table evaluation.  The lower limit is
the first abscissa of the generated table in this release. \\
\path{--help}, \path{-h} & no value & Print the same complete interface from
the executable and exit successfully; no collision strength is required. \\
\end{longtable}
\normalsize

\paragraph{Material-SIDM command line.}
Material tasks use the common restart, output, endpoint, timestep, cadence, and
radial-mesh options in the first part of the table.  They reject every
\code{--subhalo-*} option and additionally accept:
\begin{itemize}[leftmargin=1.5em]
\item \code{--sidm-collision-strength S}: required and positive.
\item \code{--cross-section-model MODEL}: required only for
\code{venus_gpu_sidm_boltzmann_poisson}; choose \code{constant},
\code{power_law}, or \code{yukawa_born}.  The other three material executables
pin this choice.
\item \code{--sidm-sigma0 X}: positive constant or power-law normalization.
\item \code{--sidm-sigma-power P}: positive power-law exponent; default 4.
\item \code{--sidm-sigma-u0 U}: positive power-law softening speed.
\item \code{--sidm-yukawa-alpha A}: positive effective kernel amplitude;
      default \(0.01\).
\item \code{--sidm-yukawa-mchi M}: positive numerical shape parameter;
      default 1.
\item \code{--sidm-yukawa-mphi M}: positive numerical shape parameter;
      \(m_\phi/m_\chi\) sets the speed scale; default 1.
\item \code{--sidm-collision-interval DT}: positive fixed collision-event
interval; default \(0.5\).
\item \code{--sidm-collision-cfl C}: per-substep probability bound in
\((0,1]\); default \(0.2\).
\item \code{--sidm-collision-maximum-substeps N}: positive safety ceiling;
default 16,384.
\item \code{--sidm-adaptive-shell-timestep}: give every fixed-occupancy
material shell an independent local CFL clock; valid for every GPU material
cross-section law.
\item \code{--sidm-target-shell-markers N}: target local shell occupancy;
default 4,096.
\item \code{--sidm-minimum-shell-markers N}: minimum shell occupancy; default
32 and no larger than the target.
\item \code{--sidm-fail-on-nonlocal-shell}: stop when a collision shell is
wider than its mean free path.
\end{itemize}
The constant and power-law tasks require \code{--sidm-sigma0}; power law also
requires \code{--sidm-sigma-u0}; Yukawa--Born uses the three amplitude/shape
parameters and rejects \code{--sidm-sigma0}.  All material tasks require a
static restart whose marker masses represent total density.

For \(F>0\), the event-clock interval represented by the two collision-cadence
options is
\begin{equation}
 \Delta t_{\rm coll}=\min\!\left[\Delta t_{\max},
 F\min(\tau_{\rm collision},\tau_{\rm dyn},\tau_{\Phi})\right],
\end{equation}
when the controlling timescale is finite; otherwise the maximum interval is
used.  For \(F=0\),
\(\Delta t_{\rm coll}=\Delta t_{\max}\).  Collisionless calculations use
\code{venus_gpu_vlasov_poisson}; the explicit subhalo and material task names
require a positive strength.
For \(S>0\), VENUS also requires the selected interval to exceed
\(128\epsilon_{\rm mach}\max(1,|t|)\).  If it does not, advancing the
floating-point time coordinate is no longer representable; the runner writes
a controlled cadence-stop checkpoint instead of entering a zero-advance
loop or silently skipping collisions.

For a continuation, repeat the original mesh, collision and cadence
arguments so the saved particle clocks retain their physical meaning.  The runner verifies the physics task, universe, Poisson source
and collision model against existing metadata.  Retain the original command
with the run so that the remaining physical parameters and numerical
controls are supplied consistently.  The radial mesh contract and metadata
remain associated with the initial segment.

\subsection{GPU input and output contract}

GPU library users configure public C++ option structures; the task executables
use command-line options. They load particle initial data from CPU-generated
weighted-PIC restarts.  The version-6 restart uses magic \code{VENUSWVPIC1}, a version and
endian sentinel, metadata clocks/counters, and one complete record per marker.
The legacy magic is intentionally unchanged so CPU and v1 checkpoints remain
interoperable.
GPU readers accept versions 1--6. Version 7 stores the CPU static
Boltzmann event solver's frozen field and evolution controls; that
integration path has a separate continuation contract.

The reader derives the version-specific marker-record size and requires the
declared count to equal the remaining payload byte-for-byte before resizing a
marker vector.  It rejects negative or non-finite clocks, invalid delayed-inner
return times, invalid marker states, duplicate IDs, non-finite physical fields,
and trailing bytes.  Positive infinity is permitted only for an event
timescale that has not yet been measured.  The checkpoint comparator treats
matching positive-infinity sentinels as equal, treats a finite/infinite pair as
a mismatch, and requires finite nonnegative tolerances.

Material and subhalo collision entry points validate model enumerations,
finite option domains, CUDA integer index bounds, adaptive-substep narrowing,
and counter-based RNG exhaustion before launching affected kernels.  MPI
marker migration and checkpoint gathering validate all byte counts and
displacements collectively against the legacy \code{int} limits of
\code{MPI_Alltoallv} and \code{MPI_Gatherv}; one rank cannot proceed after a
peer detects an invalid layout.

The task runner writes
\begin{itemize}
\item \path{profile_NNNNN.csv}: time, radius, density, radial mean and
variance, tangential second moment, and effective marker number;
\item \path{history_gpu.csv}: step, time, timestep clocks, mass, marker count,
and wall time;
\item \path{collision_history_gpu.csv}: trials, accepted events, scatter and
capture counts, invariant residuals, locality, collision clocks, and (for
material SIDM) shell-visit work;
\item \path{restart_state.bin}, \path{metadata.json}, and
\path{radial_mesh_contract.txt};
\item explicit \path{completed.ok}, \path{walltime_checkpoint.ok},
\path{locality_stopped.ok}, or \path{cadence_stopped.ok} terminal states.
\end{itemize}
The output moments are deposited after the final radial kick of the KDK
step. Density is unchanged by this kick, while the radial mean and variance
follow the completed endpoint velocities. This makes the profile and its
particle checkpoint describe the same state.
\codesource{GPU endpoint profile moments}
{\path{gpu/src/wpic_production_driver.cu}}
{Refresh the deposited moments from the completed particle state before output}
\begin{lstlisting}[language=C++]
vp.deposit_and_solve(state, time, vp_control);
\end{lstlisting}
For a continuation in the same directory, the runner checks the existing
profile and history times before opening histories for append. A diagnostic
later than the supplied checkpoint causes a clear error and leaves those
files intact. When a checkpoint lies between profile outputs, the completed
earlier profile is preserved and evolution resumes toward the next scheduled
output. An earlier checkpoint can be evolved separately in a new output
directory. The signals \code{SIGUSR1} and \code{SIGTERM} request an atomic
checkpoint followed by exit status 99.
\codesource{GPU continuation time check}
{\path{gpu/src/wpic_production_driver.cu}}
{Verify that existing diagnostics precede the supplied particle state}
\begin{lstlisting}[language=C++]
require_output_precedes_restart(run.output, source.time);
\end{lstlisting}
Plot/GIF refresh creates frame-time, last-curve-time, and input-count
sidecars.  Those sidecars define the animation's physical-time mapping.

\subsection{GPU scope limits}

The task interfaces and scientific comparison have the following scope:
\begin{itemize}
\item material MPI changes the fixed-occupancy shell populations and requires
rank-count convergence;
\item task runners use zero central point mass, origin parity and a reflecting
outer boundary. Expanding particle masses already represent density contrast;
central-source and outflow library controls require their own validation.
The delayed-return method remains in the CPU implementation;
\item long chaotic Monte Carlo comparisons use conservation, resolved profiles,
event rates and seed ensembles to test agreement of the predicted evolution;
\item Einstein--Vlasov, SphDSMC, canonical-orbit and grid collision backends
remain CPU implementations.
\end{itemize}
